\section{Discussion}
\label{sec:discussion}

This section asks three questions. What does the maize result mean, and what could make it wrong? How far does it travel beyond the 17 departments studied? And what did the way the project was run teach, independently of its results?

\subsection{What the maize result means}

In the main sample, irrigated maize keeps about \IrrShareOfRainfedSens\,\% of the heat sensitivity of rainfed maize: about $-$\QoneIrrigatedPct\,\% per degree-day above 32\,$^{\circ}$C against $-$\QoneRainfedPct\,\%. The direction matches studies abroad, which find a much smaller heat sensitivity for irrigated than for rainfed maize in the United States \parencite{schauberger2017} and, for wheat in India, about a quarter of the rainfed sensitivity \parencite{zaveri2019}. The data, crops and designs differ too much for a closer comparison of magnitudes.

Two features of the result deserve emphasis. Irrigation is associated with a smaller sensitivity to heat, not with its disappearance: irrigated maize still loses yield in hot summers. And the association rests largely on a few summers, above all 2003, which carries between a quarter and two fifths of the useful variation in heat, depending on whether weights and soil moisture are taken into account (\cref{sec:exposure,subsec:maize-inference}). This is why the strength of the evidence depends on the inference: the bootstrap by year rejects, the more conservative cluster jackknife does not. A linear effect per degree-day is a convenient summary, not a law: in 2003 the observed widening of the gap was smaller than the linear estimate predicts (\cref{sec:maize}).

\paragraph{Heat, or hot and dry summers.} In France the summers with many degree-days above 32\,$^{\circ}$C (2003, 2019, 2022) were also dry. The soil wetness index absorbs part of this, but it comes from a land-surface model whose documentation does not mention irrigation, and the result holds when it is dropped or replaced by a climatic water balance (\cref{sec:maize}), but heat and the water balance remain correlated. ``Heat'' in this monograph therefore means hot, and mostly dry, summers. The vapour pressure deficit makes this precise: in a model with both, the association moves from degree-days above 32\,$^{\circ}$C to the drying power of the air (\cref{subsec:maize-vpd}), as \textcite{lobell2013} found for American maize. The two are too correlated to be separated here, but the safest reading of the main result is that irrigated maize resists hot, dry summers better, more than that it resists heat as such. Whether irrigation helps by supplying water or also by cooling the canopy cannot be told apart with these data. Satellite studies elsewhere attribute about a third of the benefit of irrigation to cooling \parencite{zhu2022}, and heat stress computed from air temperature overstates the stress of irrigated crops \parencite{siebert2017}; with air temperature common to both modes, as here, part of what is called heat may be the absence of this cooling on rainfed fields. A check with MODIS surface temperature is consistent with this reading without settling it: across the 38 departments, surface temperature rises by about \LstSlope{}\,$^{\circ}$C per degree-day above 32\,$^{\circ}$C, a little less where more maize is irrigated (\LstSlopeQone{} at the first quartile of the irrigated share, \LstSlopeQthree{} at the third; two-sided $p = \LstInterP$), but the irrigated share also differs across departments in other ways.

\paragraph{Water restrictions.} The pre-specified interaction with crisis-level restrictions has the opposite sign to the hypothesis (\cref{sec:maize}). It is read descriptively, as planned, and it gives no support to the idea that restrictions reduced the advantage of irrigated maize in these departments.

\subsection{What could make it wrong}

The estimate is an association within departments and years. For it to measure how irrigation changes the response to heat, irrigated and rainfed land would have to differ only in ways that do not themselves vary with the heat of the summer. Several mechanisms could break this.

\paragraph{Silage.} In a hot, dry summer, some rainfed maize sown for grain may be harvested as silage instead. Abandonment is less of a concern: the SAA counts unharvested area with a zero yield, so abandoned fields stay in the grain statistics. Silage is different. If the worst rainfed fields move from grain to silage, the observed rainfed grain yield is the yield of the remaining fields and overstates what rainfed maize produced. Under that assumption, this would bias the gap downwards in hot years, against the result; if the fields moved to silage are not the worst ones, the direction of the bias is not known. The SAA areas of silage maize, available for 2010--2025, show no detectable increase in the share of rainfed maize harvested as silage in hotter summers (\cref{subsec:maize-composition}), but 2003 is outside that period and the areas may be recorded before harvest. The irrigated share of grain maize area does not move detectably with heat (\cref{sec:maize}), but areas are largely decided at sowing, so this check cannot see the mechanism.

\paragraph{Irrigation itself responds to heat.} Farmers irrigate more in hot summers, unless restrictions stop them. The estimate therefore mixes the agronomic effect of water with the behavioural response of irrigators. This is part of what ``irrigated maize'' means in practice, and it is why the result describes the observed advantage of irrigated maize, not the effect of a fixed amount of water.

\paragraph{Other changes in the departments.} The year effects absorb national trends, not changes specific to a department that would move with its summers. One such change was examined: wind farms, which an American study associates with higher nearby yields through their effect on local climate \parencite{kaffine2019}. Onshore wind capacity per hectare of farmland, built from the national register of electricity installations, shows no positive association with irrigated, rainfed or wheat yields and does not move the gap. Its one coefficient with $p < 0.05$, for irrigated maize, is negative and does not survive a correction for the four tests. The test is weak in any case: the register lists only installations in service in 2026, and departments are too coarse for a local effect.

\paragraph{Different land, different farms.} Irrigated fields tend to be on soils and farms suited to intensive maize. Department fixed effects absorb the stable part of these differences. They would not absorb, for example, a tendency of better farms to protect their irrigated maize more in hot years.

\paragraph{Measurement.} Four sources of error are documented in the monograph:
\begin{itemize}
  \item the result depends on measuring heat where maize grows (\cref{sec:maize});
  \item SIM2 is an 8~km reanalysis; E-OBS, interpolated from stations, counts more degree-days above 32\,$^{\circ}$C, and the result holds with both;
  \item the masks do not distinguish irrigated from rainfed parcels, so both modes are assigned the same heat;
  \item fixed calendar windows ignore year-to-year shifts in crop development \parencite{leroux2024}.
\end{itemize}
Errors in heat that are unrelated to the gap would bias $\beta$ towards zero, and instrumenting the SAFRAN measure by E-OBS raises the estimate by about a third (\cref{subsec:maize-measure}); errors that differ systematically between irrigated and rainfed land have no known direction.

\paragraph{The data on yields.} The quality status of 2025 is applied to 2000--2009, two releases of the statistics are spliced in 2010, many yields are rounded, and the 2025 values are provisional. The splice was checked on the overlap and the results do not depend on 2025, but the older decade rests on a release that was never revised.

\paragraph{How the split between irrigated and rainfed maize is produced.} The SAA is a synthesis of several sources, compiled by regional statistical services. This monograph does not document how they set the irrigated and rainfed yields separately, nor whether ``irrigated'' means land actually irrigated that year or land equipped for irrigation. If part of the split were estimated from expert judgement about dry summers, the gap could reflect the estimation method as much as the harvest. This was flagged when the project was scoped and remains the main untested limitation of the outcome variable.

\subsection{How far the result travels}

The main sample covers \MainShareArea\,\% of the maize area of the population, and within it the significance of the association rests on the six departments with uplands: without them it is not significant, although their slope does not differ detectably from that of the others (\cref{subsec:maize-measure}). It is concentrated in Alsace, the Rhône--Alpes and Auvergne area, the Centre and Burgundy--Franche-Comté. It excludes the South-West (Landes, Pyrénées-Atlantiques, Gers, Lot-et-Garonne, among others), France's main maize basin, and the Centre-West (Charente, Charente-Maritime, Deux-Sèvres, Vendée, Vienne), whose irrigated series have a lower quality status. With all 38 departments, the estimate is larger (\QoneAllDepsPct\,\% per degree-day) and significant, which is reassuring but rests on data the statistics office itself rates lower. Beyond grain maize, the same association appears for silage maize in \FodderDeps{} departments over 2010--2025 (\FodderPct\,\% per degree-day), so it is not specific to one crop's statistics.

The estimand also limits extrapolation in time. It describes the response to a hot summer given current varieties, practices and water rules. A lasting rise in temperature would bring adaptation, in sowing dates, varieties and the extent of irrigation, and possibly tighter water restrictions. Whether panel estimates of this kind approximate such long-run responses depends on conditions studied by \textcite{merel2021}; I do not claim that they hold here.

\subsection{What would settle the open questions}

\begin{itemize}
  \item \textbf{A map of irrigated parcels.} Assigning heat separately to irrigated and rainfed fields would remove the main measurement limitation. The parcel register used here does not record irrigation; maps derived from Sentinel-1 and Sentinel-2 exist for some basins, not for the whole country.
  \item \textbf{Surface temperature on irrigated and rainfed fields separately.} The MODIS check above averages both modes within each pixel; combined with a map of irrigated parcels, surface temperature would show directly whether irrigated fields stay cooler in hot summers \parencite{zhu2022}.
  \item \textbf{How the SAA splits irrigated and rainfed maize.} The public descriptions of the SAA do not say whether ``irrigated'' means irrigated that year or equipped for irrigation. The statistical service could answer it, and its farm-practice surveys, which record irrigation and yield field by field, could validate the split at the regional level.
  \item \textbf{Data on maize moved from grain to silage in dry years.} The SAA areas of silage maize show no shift over 2010--2025; field-level data on the harvest decision, and data covering 2003, would test the main selection threat directly.
  \item \textbf{Water volumes.} The national register of water withdrawals offers annual irrigation volumes for 2012--2023 for all water agencies, but for all crops together and with breaks in reporting; I audited this source for the project but did not use it. It could support a descriptive value per cubic metre in departments where maize dominates irrigation, not an estimate of the productivity of water.
  \item \textbf{More hot summers.} The information on heat grows only when hot summers occur. Re-estimating the model as new harvests are published, with the specification unchanged, is the cleanest test of the result.
\end{itemize}

\subsection{What the way of working taught}

Three lessons are independent of the results.

\paragraph{Check the size of your tests on your data.} The usual standard errors would have made the wheat result look solid ($p = \WheatPdep$ instead of $\WheatPwcr$). A placebo with permuted years costs a few minutes and changed the conclusions of this project. It does not detect every problem (serial correlation is one it cannot see), but it is a check that panels with common shocks should not skip. Nor should a single valid test be the end of it: when one year dominates, as 2003 does, other valid inferences can disagree, and reporting them is part of the result.

\paragraph{Choose the benchmark before the model.} In forecasting, the pre-specified trend benchmark was weaker than the plain department mean. Judged against it, a model looked useful when it was not. The comparison that matters is with the simplest forecast a user would otherwise make.

\paragraph{Pre-specification helps, within limits.} Writing the specification first, with its hash in a log, made it visible, and costly, to choose the variant that worked: the wheat result without 2025 was significant in the one-sided test and was not retained. It did not make the work blind, and it was not timestamped by anyone but the author: the specifications entered the public repository together with the first results. The maize question was chosen after an exploratory analysis, some checks were added after results were known, and one departure went unnoticed until a review. Saying so, and listing every departure (\cref{app:deviations}), is part of what makes the remaining claims credible.
